\section{Évaluation}

L'évaluation repose sur la soutenance : la compréhension du code présenté
et la capacité à répondre aux questions posées sur son fonctionnement. Un
code qui fonctionne mais n'est pas compris par son auteur n'est pas
considéré comme acquis.

\subsection{Pendant la réalisation et la soutenance}
\begin{itemize}
  \item Le code doit être écrit et compris par l'étudiant lui-même : pas d'intelligence artificielle.
  \item Pas de téléphone.
\end{itemize}

\section{Règles de codage}

Ces règles s'appliquent à tous les projets de ce dossier. Elles sont
vérifiées lors de la relecture du code et discutées en soutenance.

\subsection{Pas de variable globale non justifiée}

Toute donnée nécessaire à une fonction lui est transmise par paramètre
(ou renvoyée par valeur de retour), jamais via une variable globale. Une
variable globale ne se justifie que dans des cas rares et documentés
(exemple : une constante de configuration). Dans ce cas, le commentaire
explique pourquoi elle ne peut pas être un paramètre, et son nom est
préfixé \texttt{g\_} (exemple : \texttt{g\_nb\_joueurs\_max}). Ce préfixe
est la seule exception à la règle \og pas de préfixe \fg{} de la section
suivante : une variable globale doit sauter aux yeux partout où elle est
utilisée.

\subsection{Noms de variables, fonctions et tableaux (inspiré de la norme 42)}

Un nom doit permettre de comprendre le rôle de l'identifiant sans lire le
reste du code. Les variables locales et les tableaux ne portent pas de
préfixe de type façon notation hongroise (\texttt{l\_}, \texttt{tab\_}...) :
le rôle se lit dans le nom, pas dans un préfixe. Les seuls préfixes
autorisés sont ceux listés ci-dessous, chacun avec un sens précis.

\begin{itemize}
  \item \textbf{snake\_case partout}, jamais de camelCase : \texttt{nb\_joueurs}, pas \texttt{nbJoueurs}.
  \item \textbf{Variables} : nom explicite et complet. Préférer \texttt{nb\_joueurs} à \texttt{n}, \texttt{indice\_carte\_tiree} à \texttt{i2}. Seuls \texttt{i}, \texttt{j}, \texttt{k} sont tolérés, uniquement comme compteur de boucle local et de portée très courte (\texttt{for (int i = 0; ...)}).
  \item \textbf{Une seule déclaration par ligne} : \texttt{int a; int b;}, jamais \texttt{int a, b;}.
  \item \textbf{Booléens} : préfixés par \texttt{est\_}, \texttt{a\_} ou \texttt{peut\_}, pour que leur lecture soit une question ou une affirmation : \texttt{est\_gagnant}, \texttt{a\_rejoue}, \texttt{peut\_tirer}.
  \item \textbf{Tableaux} : nom explicite, au pluriel si possible, décrivant ce qu'il contient, sans préfixe \texttt{tab\_} : \texttt{scores}, \texttt{cartes\_tirees}, \texttt{grille}.
  \item \textbf{Fonctions} : \texttt{verbe\_nom}, une action claire par fonction : \texttt{calculer\_score()}, \texttt{afficher\_plateau()}. Une fonction qui renvoie un booléen se lit comme une question : \texttt{est\_victoire()}, \texttt{peut\_jouer()}.
  \item \textbf{Constantes et macros} : \texttt{UPPER\_SNAKE\_CASE} : \texttt{NB\_JOUEURS\_MAX}, \texttt{TAILLE\_GRILLE}.
  \item \textbf{Variable globale justifiée} : préfixe \texttt{g\_} (voir ci-dessus).
  \item \textbf{\texttt{struct}, \texttt{enum}, \texttt{union}, \texttt{typedef}} : préfixe \texttt{s\_} pour une struct, \texttt{e\_} pour un enum, \texttt{u\_} pour une union, \texttt{t\_} pour le nom créé par \texttt{typedef} (exemple : \texttt{typedef struct s\_joueur \{ ... \} t\_joueur;}). Ces constructions ne sont pas obligatoires dans ces projets ; si elles sont utilisées, ces préfixes s'appliquent.
\end{itemize}

\subsection{Indentation}

Chaque niveau d'imbrication (fonction, boucle, condition) est indenté de
façon cohérente et visible, de sorte que la structure du code soit
lisible sans compter les accolades. Un éditeur de code réindente
généralement un fichier automatiquement ; prendre l'habitude de le faire
avant de rendre son travail.

\begin{UPSTIinfor}{Pour aller plus loin : éviter les nombres magiques}
Une valeur qui a un sens précis dans le programme (une taille de grille,
un code de carte, un nombre maximal de joueurs...) est plus lisible et
plus facile à modifier si elle porte un nom, plutôt que d'apparaître en
dur à plusieurs endroits du code. La directive \texttt{\#define} n'ayant
pas encore été vue en cours au moment de ces projets, ce point est une
information, pas une règle vérifiée en soutenance.
\begin{lstlisting}[language=C]
#define TAILLE_GRILLE 4

int grille[TAILLE_GRILLE][TAILLE_GRILLE];
\end{lstlisting}
\texttt{TAILLE\_GRILLE} se comporte comme la valeur \texttt{4} partout où
elle apparaît, mais en changeant une seule ligne, toute la grille change
de taille.
\end{UPSTIinfor}

\section{Rendre son projet}

À la fin du projet, le code est rendu avec la commande \texttt{submit50}, comme pour les TP. Sur la page GitHub du projet, cliquer sur le bouton \texttt{submit50} (à côté du bouton \texttt{wget} qui a servi à démarrer), puis coller la commande dans le terminal (\texttt{Maj-INSER}) depuis le dossier du projet. Appuyer sur Entrée.

\begin{itemize}
  \item Un seul projet par rendu : se placer dans le dossier du projet choisi avant de lancer la commande.
  \item Le rendu peut être refait autant de fois que nécessaire : c'est la \textbf{dernière} version soumise avant la soutenance qui est consultée.
  \item Vérifier que le code \textbf{compile} avant de le soumettre.
  \item Soumettre son projet ne dispense pas de la soutenance : il faut pouvoir expliquer son code.
\end{itemize}
